% ============================================================
% 1. Fundamentos de la computación paralela (Sánchez)
% ============================================================
\section{Fundamentos de la computación paralela}

\begin{frame}{¿Qué es la computación paralela?}
  \begin{itemize}
    \item Divide un problema en \textbf{subtareas} que se ejecutan
          \textbf{al mismo tiempo} en varios procesadores o núcleos.
    \item Objetivo: \textbf{reducir el tiempo de cómputo} y aprovechar mejor
          los recursos.
    \item Exige \textbf{coordinación} entre quienes ejecutan las subtareas.
  \end{itemize}

  \vspace{0.3cm}
  \begin{columns}[T]
    \begin{column}{0.31\textwidth}
      \begin{block}{Memoria compartida}
        \small Núcleos de un mismo equipo comparten la memoria.
      \end{block}
    \end{column}
    \begin{column}{0.31\textwidth}
      \begin{block}{Memoria distribuida}
        \small Cada nodo tiene su memoria y se comunica por mensajes.
      \end{block}
    \end{column}
    \begin{column}{0.31\textwidth}
      \begin{block}{Híbrida}
        \small Combina ambas, como en supercomputadoras y clústeres HPC.
      \end{block}
    \end{column}
  \end{columns}

  \vspace{0.2cm}
  {\tiny Fuente: \cite{dai2025parallel}, sec.~2.1.}
\end{frame}

\begin{frame}{Paralelismo, concurrencia y sistemas distribuidos}
  \small
  \begin{tabularx}{\textwidth}{@{}p{2.6cm}XX@{}}
    \toprule
    \textbf{Concepto} & \textbf{Idea central} & \textbf{En una tienda en línea} \\
    \midrule
    Concurrencia &
      Varias tareas avanzan en un mismo período; pueden alternarse
      \cite{reinders2023}. &
      Un servidor atiende varios pedidos intercalando su ejecución. \\
    Paralelismo &
      Varias tareas se ejecutan en el mismo instante en recursos distintos
      \cite{reinders2023,dai2025parallel}. &
      Tres trabajadores procesan tres pedidos a la vez. \\
    \raggedright Sistema distribuido &
      Elementos autónomos que se comunican por red y se perciben como un
      solo sistema \cite{vansteen2023distributed}. &
      Tienda, broker y trabajadores en máquinas distintas. \\
    \bottomrule
  \end{tabularx}

  \vspace{0.3cm}
  \begin{block}{Relación}
    Se complementan: un sistema distribuido puede ejecutar en paralelo dentro
    de cada nodo; el paralelo prioriza velocidad y el distribuido,
    escalabilidad y tolerancia a fallos \cite{dai2025parallel}.
  \end{block}
\end{frame}

\begin{frame}{Ejemplo: varios trabajadores procesando pedidos}
  \begin{columns}[T]
    \begin{column}{0.47\textwidth}
      \textbf{Un trabajador (secuencial)}\\[0.15cm]
      \small
      \begin{tabular}{@{}ll@{}}
        \toprule
        Turno & Pedido \\
        \midrule
        1 a 6 & P1 $\rightarrow$ P2 $\rightarrow$ \dots $\rightarrow$ P6 \\
        \bottomrule
      \end{tabular}\\[0.15cm]
      Total: \textbf{6 turnos}.
    \end{column}
    \begin{column}{0.47\textwidth}
      \textbf{Tres trabajadores (paralelo)}\\[0.15cm]
      \small
      \begin{tabular}{@{}lll@{}}
        \toprule
        Trabajador & Turno 1 & Turno 2 \\
        \midrule
        T1 & P1 & P4 \\
        T2 & P2 & P5 \\
        T3 & P3 & P6 \\
        \bottomrule
      \end{tabular}\\[0.15cm]
      Total: \textbf{2 turnos}.
    \end{column}
  \end{columns}

  \vspace{0.3cm}
  \begin{block}{En RabbitMQ}
    Los trabajadores son consumidores de una \textbf{misma cola}: el broker
    reparte los mensajes entre ellos y entrega cada pedido a un solo
    trabajador a la vez \cite{rabbitmq2026amqp}.
  \end{block}

  {\tiny Ejemplo ilustrativo: supone pedidos independientes y de igual
  duración; no considera el costo de coordinación.}
\end{frame}

\begin{frame}{Ejemplo: pedidos de distinta duración}
  \small
  P1 y P4 tardan \textbf{4 min}; P2, P3, P5 y P6, \textbf{1 min}.
  Trabajo total: 12 min.\\[0.2cm]

  \begin{tabularx}{\textwidth}{@{}p{3.7cm}XXXc@{}}
    \toprule
    \textbf{Reparto} & \textbf{T1} & \textbf{T2} & \textbf{T3} &
      \textbf{Termina} \\
    \midrule
    \raggedright Un solo trabajador & \multicolumn{3}{l}{P1 a P6} & 12 min \\
    \raggedright Por turnos (predeterminado) & P1, P4 (8) & P2, P5 (2) &
      P3, P6 (2) & 8 min \\
    \raggedright Al quedar libre (\texttt{prefetch} = 1) & P1 (4) &
      P2, P4 (5) & P3, P5, P6 (3) & \textbf{5 min} \\
    \bottomrule
  \end{tabularx}

  \vspace{0.3cm}
  \begin{block}{Qué muestra}
    Tres trabajadores no bastan para terminar en un tercio del tiempo: importa
    cómo se reparten los pedidos. Por defecto, RabbitMQ entrega por turnos sin
    mirar cuánto trabajo tiene cada uno; con \texttt{prefetch} = 1 solo le
    entrega otro pedido a un trabajador después de su \texttt{ack}
    \cite{rabbitmq2026workqueues}.
  \end{block}

  {\tiny Ejemplo propio con duraciones supuestas; entre paréntesis, minutos
  ocupados por cada trabajador. El mínimo posible sería 4 min.}
\end{frame}
